\documentclass[11pt]{article}
\usepackage[margin=1in]{geometry}
\usepackage{amsmath,amssymb,amsthm,microtype,booktabs}
\usepackage[hidelinks]{hyperref}
\allowdisplaybreaks[2]

\newtheorem{theorem}{Theorem}[section]
\newtheorem{lemma}[theorem]{Lemma}
\newtheorem{corollary}[theorem]{Corollary}
\newcommand{\e}{\mathrm e}
\newcommand{\ii}{\mathrm i}
\newcommand{\dd}{\mathrm d}
\newcommand{\R}{\mathbb R}
\newcommand{\arcosh}{\operatorname{arcosh}}

\title{Volterra dynamics and logarithmic coefficient bounds\\
       for the Newman heat flow}
\author{Edward Baker}
\date{14 November 2026}

\begin{document}
\maketitle

\begin{abstract}
Orthogonalization for the complete theta source at each Newman time
gives a moving sequence of Jacobi recurrence coefficients. We derive
its Volterra equation and entrywise Lax identity. A bound on the
logarithmic derivative of the complete weight, combined with monic
polynomial integration by parts and Jensen's inequality, yields
\(\beta_n+\beta_{n+1}\le\arcosh^2(2n+12)/16\) for all \(n\ge0\) and
\(0\le\lambda\le1/2\). The resulting monic norm product gives an
explicit factor in the Gaussian quadrature error for the fixed
initial Jacobi chain. Positive weighted-square sources with
nondegenerate collisions satisfy the same coefficient bound.
\end{abstract}

\section{Introduction}

The de Bruijn--Newman field is the Fourier transform of a complete
theta density tilted by \(\e^{\lambda u^2}\). Encoding multiplication by
the source variable \(u\) in an orthogonal polynomial basis produces a
tridiagonal Jacobi matrix. One can hold the basis fixed and put a heat
factor in the readout, or rebuild the basis at each time and study the
resulting coefficient evolution. The moving basis gives the classical
Volterra lattice, the quadratic even-weight specialization of the
Toda hierarchy \cite[\S2.3]{vanassche}. The lattice describes how the
coefficients change in time. To bound them in degree, we use the
confinement supplied by the complete theta density.

The key estimate bounds a sum of adjacent squared off-diagonal
coefficients by an explicit function growing as the square of a
logarithm. It follows from a scalar identity for the probability measure
whose density is the weight times the square of a monic polynomial,
normalized by its squared norm. The estimate controls a norm factor in
a real-height quadrature error bound. Its proof also applies to a class
of positive weighted-square densities, including examples whose fields
collide at positive heat times.

\section{The complete source and its polynomial basis}

For \(u\ge0\), put \(\rho(u)=\pi\e^{4u}\), \(r_m=m^2\rho(u)\), and
\begin{equation}
 s_m(u)=\e^u(2r_m^2-3r_m)\e^{-r_m},\qquad
 \Phi(u)=\sum_{m\ge1}s_m(u),\qquad W_\lambda(u)=\e^{\lambda u^2}\Phi(u).
 \label{eq:source}
\end{equation}
The complete source has a smooth even extension. Indeed, for
\(\vartheta(t)=\sum_{m\in\mathbb Z}\e^{-\pi m^2t}\) let
\(p(u)=\e^u(\vartheta(\e^{4u})-1)/2\).
Jacobi's inversion formula \cite[Eqs.~20.2.3 and 20.7.32]{dlmf}
gives \(p(-u)-p(u)=\sinh u\) and hence \(p'(0)=-1/2\).
Termwise differentiation on compact real intervals gives
\(\Phi=(p''-p)/8\), proving evenness. On \(u\ge0\) all
summands are strictly positive since \(r_m\ge\pi>3\).
In particular the physical field and the reversed clock are
\begin{equation}
 H_\lambda(z)=\frac12\int_\R W_\lambda(u)\e^{\ii zu}\,\dd u
      =\int_0^\infty W_\lambda(u)\cos(zu)\,\dd u,\qquad
 \tau=\frac12-\lambda .
 \label{eq:field}
\end{equation}
The full-line mass \(Z_\lambda=\int_\R W_\lambda(u)\,\dd u\) is retained in
the physical field.

To check the time-zero scale, set \(I(z)=\int_0^\infty p(u)\cos(zu)\,\dd u\)
and \(s=(1-\ii z)/2\). Two integrations by parts give
\(8H_0(z)=1/2-(1+z^2)I(z)\), using the endpoint \(p'(0)=-1/2\).
For \(\operatorname{Re}s>1\), termwise Mellin integration gives
\(\pi^{-s/2}\Gamma(s/2)\zeta(s)\) as the integral of
\((\vartheta(t)-1)t^{s/2-1}/2\) on \(t>0\). Splitting at \(1\)
and applying theta inversion below \(1\) gives
\[
 \pi^{-s/2}\Gamma(s/2)\zeta(s)=\frac{1}{s(s-1)}
  +\int_1^\infty\frac{\vartheta(t)-1}{2}
       (t^{s/2-1}+t^{(1-s)/2-1})\,\dd t .
\]
The integral is \(8I(z)\) by \(t=\e^{4u}\) and is entire in \(s\).
Multiplication by \(s(s-1)/2\) gives the continued completed function
\(\xi(s)=1/2-(1+z^2)I(z)\), so \(H_0(z)=\xi((1-\ii z)/2)/8\).
The positive-half-line integrals and the boundary terms at infinity
are controlled by double-exponential decay.

The real-variable bounds needed below follow directly from the sum.
For \(j=0,1\) the \(j\)-th derivative of \(s_m\) is \(\e^u\)
times a polynomial of degree \(2+j\) in \(r_m\) times \(\e^{-r_m}\).
In fact the polynomials start with \(2r^2-3r\), and differentiation
sends a polynomial \(Q\) to \((1-4r)Q+4rQ'\). For \(u\ge0\) the
absolute value of each differentiated term is bounded by a constant
times
\[
 \e^u\rho(u)^{2+j}\e^{-\rho(u)}
     m^{4+2j}\e^{-\pi(m^2-1)}.
\]
The \(m\)-factor is summable. This gives local uniform convergence
and differentiation, and the same majorant argument works on any
compact real interval before reflection. Writing
\(\rho(u)=\pi\e^{4|u|}\) on the full line, we obtain
\begin{equation}
 |\Phi^{(j)}(u)|\le C_j\,\e^{|u|}\rho(u)^{2+j}\e^{-\rho(u)}
       \qquad (j=0,1)
 \label{eq:tail}
\end{equation}
with constants \(C_j\) independent of \(u\). In particular for every
fixed bound on \(|\lambda|\), any fixed polynomial times \(W_\lambda\)
or its \(u\)-derivative tends to zero at both ends and is integrable,
even after multiplication by \(\e^{b|u|}\) for fixed \(b\).
The reason is that \(-\rho(u)\) dominates the logarithm of the other
factors, which grows at most quadratically in \(|u|\).
These bounds justify differentiating the moments in \(\lambda\) on
compact sets. They also give
\(\partial_\lambda H_\lambda=-\partial_z^2H_\lambda\).

For real \(\lambda\), let \(P_n(u;\lambda)\) be the monic
orthogonal polynomial of degree \(n\) for \(W_\lambda(u)\,\dd u\).
All moment matrices are positive definite by strict positivity
of the weight and integrability of polynomials squared. This
gives existence and uniqueness at every finite degree. Set
\[
 h_n(\lambda)=\int_\R P_n(u;\lambda)^2W_\lambda(u)\,\dd u,\qquad
 \beta_n(\lambda)=\frac{h_n(\lambda)}{h_{n-1}(\lambda)}\ (n\ge1),
 \qquad \beta_0=0 .
\]
Evenness and uniqueness give \(P_n(-u;\lambda)=(-1)^nP_n(u;\lambda)\).
Expanding \(uP_n\) in the monic basis gives
\begin{equation}
 uP_n=P_{n+1}+\beta_nP_{n-1}
 \label{eq:recurrence}
\end{equation}
with the second term omitted at \(n=0\). In the expansion the coefficient
of \(P_n\) vanishes by parity, and degrees \(j<n-1\) vanish because
\(uP_j\) has degree less than \(n\). The coefficient of \(P_{n-1}\)
is \(h_n/h_{n-1}\) by testing against \(P_{n-1}\) and moving
multiplication to the other factor. With \(p_n=P_n/\sqrt{h_n}\),
the Jacobi matrix \(J(\lambda)\) in the orthonormal basis has
zero diagonal and edge between \(n-1\) and \(n\) equal to
\(a_n=\sqrt{\beta_n}\); set \(a_0=0\).

\section{The moving-basis evolution}

The coefficients of each \(P_n\) solve a finite Gram linear system:
write \(P_n=u^n+\sum_{j<n}c_j u^j\) and impose orthogonality to
\(1,u,\ldots,u^{n-1}\). Moment differentiability and positivity of
the Gram determinant give differentiability of all \(c_j\) for real
\(\lambda\). In particular all derivatives used at the physical time
endpoints are available by extending \(\lambda\) locally.

For the monic derivative, differentiate
\(\langle P_n,P_j\rangle_\lambda=0\) at \(j<n\), with
\(\langle\cdot,\cdot\rangle_\lambda\) the real weighted inner product.
The term containing \(\partial_\lambda P_j\) vanishes by degree. Thus
\[
 \langle\partial_\lambda P_n,P_j\rangle_\lambda
    =-\langle P_n,u^2 P_j\rangle_\lambda.
\]
For \(j<n-2\) the second factor has degree less than \(n\).
For \(j=n-1\) parity gives zero, and for \(j=n-2\) the second
factor has degree \(n\) and leading coefficient one. This gives
\begin{equation}
 \partial_\lambda P_n=-\beta_n\beta_{n-1}P_{n-2}\quad(n\ge2),
 \qquad \partial_\lambda P_0=\partial_\lambda P_1=0 .
 \label{eq:monicvariation}
\end{equation}
When differentiating \(h_n\) the term with \(\partial_\lambda P_n\)
also vanishes by degree. Squaring \eqref{eq:recurrence} in the
weighted norm yields
\begin{equation}
 \partial_\lambda\log h_n
  =\frac{\int u^2P_n^2W_\lambda\,\dd u}{h_n}
  =\beta_n+\beta_{n+1}\qquad(n\ge0).
 \label{eq:hflow}
\end{equation}
The term with \(P_{n-1}\) has squared norm \(\beta_n h_n\)
for \(n\ge1\). Subtracting adjacent identities in \eqref{eq:hflow}
and differentiating \(a_n=\sqrt{\beta_n}\) give
\begin{equation}
 \partial_\lambda\beta_n=\beta_n(\beta_{n+1}-\beta_{n-1}),\qquad
 \partial_\lambda a_n=\frac{a_n}{2}(\beta_{n+1}-\beta_{n-1})
 \qquad(n\ge1).
 \label{eq:volterra}
\end{equation}
In the reversed time \(\tau\), both right sides change sign.

The Lax relation follows directly from the same basis derivative.
Write \(J^2=L+D+U\) with strictly lower part \(L\), diagonal \(D\)
and strictly upper part \(U\). Then
\(D_{nn}=\beta_n+\beta_{n+1}\) and
\(L_{n,n-2}=a_n a_{n-1}\) for \(n\ge2\), with all other strictly
lower entries zero. By \eqref{eq:monicvariation} and \eqref{eq:hflow},
the polynomial column \(\boldsymbol p=(p_n)_{n\ge0}\) satisfies
\[
 \partial_\lambda\boldsymbol p=A\boldsymbol p,\qquad
 A=-L-D/2=B-J^2/2,\qquad B=(U-L)/2 .
\]
Differentiate \(J\boldsymbol p=u\boldsymbol p\) row by row.
Using \(u A\boldsymbol p=A J\boldsymbol p\) gives
\((J'-[A,J])\boldsymbol p=0\).
Independence of the polynomials in every finite row, and
\([J^2,J]=0\), yield
\begin{equation}
 \partial_\lambda J=[B,J],\qquad [B,J]=BJ-JB .
 \label{eq:lax}
\end{equation}
The matrix \(B\) is skew-symmetric. The two banded products in the
commutator act on finite sequences and every entry uses finitely
many coefficients. This domain suffices for \eqref{eq:lax}; a
time-dependent bounded-operator or unitary evolution would require
further domain and regularity analysis.

For the scalar interpretation, write \(d\mu_\lambda=W_\lambda\,\dd u/Z_\lambda\).
Polynomials are dense in \(L^2(\mu_\lambda)\). Indeed, a function
\(g\) orthogonal to all polynomials has a finite measure
\(\overline g\,\dd\mu_\lambda\) with entire Fourier transform:
Cauchy--Schwarz and \eqref{eq:tail} allow arbitrary fixed exponential
weights and derivatives. All derivatives at \(0\) vanish. The
transform is therefore zero, and uniqueness of the Fourier transform
of a finite measure implies \(g=0\) almost everywhere. For this last
step, convolution with a Gaussian gives zero by Fourier inversion,
and Gaussians approaching a point mass then recover the measure.
Multiplication by \(u\), with domain
\(\{f:uf\in L^2(\mu_\lambda)\}\), transports to a self-adjoint
realization of \(J(\lambda)\) via the orthonormal basis
\(\sqrt{Z_\lambda}p_n\). This basis maps \(1\) to \(e_0\).
In particular, for real \(x\),
\begin{equation}
 \begin{aligned}
 H_\lambda(x)&=\frac{Z_\lambda}{2}
       \langle e_0,\e^{\ii xJ(\lambda)} e_0\rangle,\\
 H_{1/2-\tau}(x)&=\frac{Z_*}{2}
       \langle e_0,\e^{-\tau J_*^2}\e^{\ii xJ_*} e_0\rangle,
       \qquad J_*=J(1/2),\quad Z_*=Z_{1/2}.
 \end{aligned}
 \label{eq:twochains}
\end{equation}
The first line uses the moving weight, while the second keeps the
initial basis and uses the factor \(\e^{-\tau u^2}\).
Both expressions are justified by multiplication in the respective
weighted spaces.

\section{Complete-source confinement and the coefficient envelope}

Set \(V_\lambda=-\log W_\lambda\).
On \(u\ge0\), the individual score is
\[
 (\log s_m)'=1+4+\frac{8r_m}{2r_m-3}-4r_m
            =9+\frac{12}{2r_m-3}-4r_m .
\]
The derivative of the complete sum divided by that sum is the
weighted mean of the summand scores, with positive weights
\(s_m/\Phi\). The estimates in \eqref{eq:tail} justify this mean.
Since \(r_m\ge\pi\e^{4u}>3\) and \(2\lambda u\le u\) on the
physical time range, we have
\begin{equation}
 V_\lambda'(u)\ge 4\pi\e^{4u}-13-u \qquad(u\ge0).
 \label{eq:score}
\end{equation}
The quantity \(uV_\lambda'(u)\) is even. With \(t=|u|\le1/2\),
\eqref{eq:score} gives \(uV_\lambda'(u)\ge-13t-t^2\ge-27/4\).
Also \(\cosh(2)<4\): its series begins \(1+2\), and the remaining
terms start at \(2/3\) with ratios bounded by \(2/15\).
Hence \(uV_\lambda'(u)\ge\cosh(4u)-11\) in this region.
For \(t\ge1/2\), \(4\pi-1/t>10\) and
\(10\e^{4t}\ge10+40t\ge13+t\). It follows that
\((4\pi t-1)\e^{4t}\ge t(13+t)\), so
\(uV_\lambda'(u)\ge\e^{4t}\ge\cosh(4t)\).
Together these estimates give
\begin{equation}
 uV_\lambda'(u)\ge\cosh(4u)-11,\qquad
 u\in\R,\quad 0\le\lambda\le1/2 .
 \label{eq:confinement}
\end{equation}

\begin{lemma}[Confinement and monic norms]
\label{lem:general}
Let \(w\) be positive, even and locally absolutely continuous on
\(\R\). Assume that for every integer \(k\ge0\),
\[
 (1+|u|)^k\bigl(w(u)+|w'(u)|+\cosh(\kappa u)w(u)\bigr)
 \quad\text{is integrable},\qquad
 (1+|u|)^k w(u)\longrightarrow 0
 \quad(|u|\to\infty)
\]
for a fixed \(\kappa>0\). If \(V=-\log w\) obeys
\(uV'(u)\ge\cosh(\kappa u)-C\) almost everywhere, its monic
recurrence coefficients satisfy
\[
 \beta_n+\beta_{n+1}
   \le\kappa^{-2}\arcosh^2(2n+1+C)\qquad(n\ge0)
\]
with the nonnegative real inverse. The assumptions imply \(2n+1+C\ge1\).
\end{lemma}

\begin{proof}
The monic recurrence follows as in \eqref{eq:recurrence}.
The measure \(d\nu_n=P_n^2 w\,\dd u/h_n\) is a probability.
Integrating \((uP_n^2w)'\) over \(\R\) gives zero by the
boundary hypothesis. Its \(w'\)-term equals
\(-uV'P_n^2w\) and is integrable by assumption. The polynomial
\(uP_n'-nP_n\) has degree less than \(n\), using a \(u\)-derivative
here. Orthogonality gives the first equality below, and
the squared recurrence norm gives the second:
\[
 \mathbb E_{\nu_n}[uV']=2n+1,\qquad
 \mathbb E_{\nu_n}[u^2]=\beta_n+\beta_{n+1}.
\]
The function \(v\mapsto\cosh(\kappa\sqrt v)\) is convex on
\(v\ge0\), as its power series consists of nonnegative multiples
of integer powers \(v^j\). Jensen and the assumed pointwise
lower bound yield
\[
 \cosh\bigl(\kappa\sqrt{\beta_n+\beta_{n+1}}\bigr)
  \le\mathbb E_{\nu_n}[\cosh(\kappa u)]\le 2n+1+C .
\]
Inversion and squaring prove the claim.
\end{proof}

\begin{theorem}[Theta coefficient bound]
\label{thm:theta}
Define \(B_n=\arcosh^2(2n+12)/16\). For the complete theta weight
and every \(0\le\lambda\le1/2\),
\[
 \beta_n(\lambda)+\beta_{n+1}(\lambda)\le B_n\qquad(n\ge0).
\]
In particular, for \(N\ge1\),
\begin{equation}
 \frac{h_N(1/2)}{h_0(1/2)}
 \le\prod_{j=0}^{N-1}B_j
 \le\left[\frac{\arcosh(2N+10)}{4}\right]^{2N}.
 \label{eq:product}
\end{equation}
\end{theorem}

\begin{proof}
In Lemma~\ref{lem:general} take \(\kappa=4\) and \(C=11\) using
\eqref{eq:confinement}. All tail hypotheses follow from
\eqref{eq:tail}, including that for the factor \(\cosh(4u)\).
Each integration by parts is at a fixed finite degree; the
constant in \eqref{eq:confinement} is common to all degrees and times.
For the product, \(\beta_j\le B_{j-1}\) by nonnegativity.
Multiply for \(1\le j\le N\) at \(\lambda=1/2\) and telescope
\(h_j/h_{j-1}\). The second inequality uses monotonicity of \(B_j\).
\end{proof}

Since \(0\le\partial_\lambda\log h_n\le B_n\) by \eqref{eq:hflow},
for \(0\le\lambda_1\le\lambda_2\le1/2\) we also have
\[
 0\le\log\frac{h_n(\lambda_2)}{h_n(\lambda_1)}
     \le(\lambda_2-\lambda_1)B_n.
\]
This is control of the norms along the specified weight deformation.

\section{A scalar consequence for the fixed chain}

Let \(J_N\) be the leading \(N\times N\) block of \(J_*\) for
\(N\ge1\), and set
\[
 H_N(\tau,x)=\frac{Z_*}{2}
       \langle e_0,\e^{-\tau J_N^2}\e^{\ii xJ_N}e_0\rangle,
 \qquad
 S_N(\tau,x)=\sum_{k=0}^N
       \frac{\tau^k x^{2N-2k}}{k!(2N-2k)!}.
\]
Here \(x\) is real and \(0\le\tau\le1/2\), with zeroth powers
interpreted as one. The normalized monic squared norm for
\(\mu_*=\mu_{1/2}\) is \(h_N(1/2)/h_0(1/2)\), since
\(P_0=1\). This is the norm that enters the quadrature bound.

\begin{corollary}
\label{cor:scalar}
For the stated real-height and time domain,
\[
 |H_{1/2-\tau}(x)-H_N(\tau,x)|
 \le\frac{Z_*}{2}\frac{h_N(1/2)}{h_0(1/2)}S_N(\tau,x)
 \le\frac{Z_*}{2}S_N(\tau,x)\prod_{j=0}^{N-1}B_j .
\]
\end{corollary}

\begin{proof}
The characteristic polynomials of the principal blocks obey
\eqref{eq:recurrence}, by tridiagonal determinant expansion, so
the eigenvalues of \(J_N\) are the zeros of \(P_N(\cdot;1/2)\).
For \(N>1\) all internal edges are positive, so the site equations
force an eigenvector with vanishing first entry to be zero.
This implies simple eigenvalues and strictly positive scalar weights
at \(e_0\); their sum is one. The same claims hold for \(N=1\),
where \(J_1=0\). Conjugating by the diagonal alternating sign gives
\(-J_N\) and fixes \(e_0\), so the scalar measure \(Q_N\) is even.
For degrees \(0\le j\le2N-1\), its moments agree with \(\mu_*\):
an expansion of the matrix moment as a sum over tridiagonal walks
shows that reaching site \(N\) and returning to \(0\) takes at least
\(2N\) steps.

For real \(s\), Hermite interpolation of \(\cos(su)\) and its first
derivative at the \(N\) nodes gives a unique polynomial \(I_s\)
of degree at most \(2N-1\), since homogeneous interpolation
forces divisibility by \(P_N^2\). Its pointwise error satisfies
\[
 |\cos(su)-I_s(u)|
     \le P_N(u;1/2)^2\,\frac{|s|^{2N}}{(2N)!}.
\]
Indeed, away from the nodes subtract
\(cP_N^2\) from the interpolation error, choosing \(c\) to produce
one additional zero at \(u\). Repeated Rolle on the interval
containing the nodes and \(u\) gives a zero of the \(2N\)-th
derivative there, so \(|c|\le |s|^{2N}/(2N)!\).
At a node the error is zero. Quadrature exactness for \(I_s\)
and agreement with \(\cos(su)\) at the nodes reduce the integral
error to the integral of this remainder against \(\mu_*\).

Finally \(\e^{-\tau u^2}\cos(xu)=\mathbb E\cos((x+\sqrt{2\tau}G)u)\)
for a standard normal \(G\). Boundedness of the cosine permits
averaging under both probability measures. Gaussian integration
by parts gives \(\mathbb E G^{2k}=(2k)!/(2^k k!)\), while the odd
moments vanish. Hence \(\mathbb E|x+\sqrt{2\tau}G|^{2N}/(2N)!
=S_N(\tau,x)\). Multiply by \(Z_*/2\) and use
\eqref{eq:product}.
\end{proof}

\section{Applicability to colliding sources}

For a finite sum \(\sum_{m=1}^M c_m s_m\) with positive \(c_m\),
multiply by \(\e^{\lambda u^2}\) on \(u\ge0\) and take the even
extension. Each constant \(c_m\) cancels in its summand score, so
the positive mean again gives \eqref{eq:score}. The extended density
is positive and locally absolutely continuous. Its tails satisfy
the hypotheses of Lemma~\ref{lem:general}. One can integrate by
parts on the two half-lines, since \(uP_n^2w\) vanishes at zero.
Thus it has the same bound as Theorem~\ref{thm:theta}.

The positive three-square and four-square sources \cite{three,four}
have validated double events in their cosine transforms at
\(\lambda=1/4\), \(44.1303<x<44.1304\), and at
\(\lambda=1/100\), \(83.5323729<x<83.5323730\), respectively.
Their Jacobi coefficients satisfy the envelope at those events.
They retain the modewise confinement used here, while the canonical
source also has equal coefficients and full modular reflection.

\section{Outstanding challenges}

\paragraph{Asymptotics in degree.}
The all-degree constant controls a sum of two adjacent coefficients.
Determining leading behavior of individual \(a_n\), with lower
bounds and quantitative remainders uniform in heat time, would
require sharper control than the confinement inequality used here.
A solution would also determine how conservative the norm product is
in quadrature for the actual theta data.

\paragraph{Stability and growing height.}
The entrywise flow controls derivatives at each finite index.
Quantitative propagation of measured coefficient errors through
growing cutoffs calls for a stability estimate that pays for the
unobserved coefficients at the boundary of a finite prefix.
The scalar quadrature bound applies for each \(N\) and real height.
Determining sharp depth and certified coefficient precision requirements
at growing heights involves the norm factors, the oscillatory readout
and the scale of the requested error.

\begin{samepage}
\paragraph{Arithmetic input and sign.}
The control examples show that the shared score inequality and
positive recurrence coefficients permit a collision. A signed
field estimate would need a source-distinguishing input and a
quantitative transfer to the relevant scalar margin. For a
logarithmic or ratio observation, the transfer would also have
to account for small denominators and zeros.
\par
\end{samepage}

\paragraph{Acknowledgment.}
Substantial OpenAI Codex assistance was dispatched as GPT-6.1 Sol /
Ultra. Separate backend attestation is unavailable.

\begin{thebibliography}{9}
\raggedright
\bibitem{dlmf}
NIST Digital Library of Mathematical Functions,
\href{https://dlmf.nist.gov/20.2.E3}{Eq.~20.2.3} and
\href{https://dlmf.nist.gov/20.7.E32}{Eq.~20.7.32}.

\bibitem{vanassche}
W.~Van Assche, \emph{Orthogonal polynomials, Toda lattices and Painlev\'e
equations}, \href{https://arxiv.org/abs/2202.11017v1}{arXiv:2202.11017v1},
2022.

\bibitem{three}
E.~Baker, \emph{Square-support constraints admit a collision even with the
first source endpoint}, Investigation~22, research note, 15 October 2026.

\bibitem{four}
E.~Baker, \emph{Two odd source endpoints still permit a positive weighted-square
collision}, Investigation~22, research note, 15 October 2026.
\end{thebibliography}
\end{document}
